% !TEX root = ../Thesis.tex

\chapter{Risultati sperimentali}
\label{ch:evaluation}

Le peculiarità del modello proposto, se paragonato con i più recenti lavori affini, risiedono primariamente nella inesigenza di questo ad addestrare un modello steganografico ad hoc, atto che porta con sé il rischio di alterare la distribuzione di probabilità delle architetture basali esistenti rendendosi quindi vulnerabile a tecniche di steganalisi, e secondariamente nell'utilizzo di un'architettura Image-to-Image quale è VQ-GAN per l'implementazione del generatore steganografico. Questa seconda peculiarità è risultata uno svantaggio a livello implementativo in quanto, sebbene il modello nel suo insieme possa essere inteso quale SwE (Steganography without Embedding), utilizzi in realtà un dataset.

This is where you show that the novel `thing' you described in Chapter~\autoref{ch:proposta} and implemented in Chapter~\autoref{ch:implementazione} is, indeed, much better than the existing versions of the same.

You will probably use figures (try to use a high-resolution version), graphs, tables, and so on. An example is shown in Figure~\ref{fig:dave}.

\begin{figure}[htbp]
\begin{center}
\includegraphics[width=.7\textwidth]{DaveSecurity}
\caption{Network Security - the sad truth}
\label{fig:dave}
\end{center}
\end{figure}

Note that, likewise tables and listings, you shall not worry about where the figures are placed. Moreover, you should not add the file extension (LaTeX will pick the `best' one for you) or the figure path.